\section{Erd\H{o}s--Hajnal reciprocal cycle-length sum (Problem 65) --- Round 10}

\subsection*{1) ROUND-10 OBJECTIVE}

\textbf{Path chosen: (A) proof-oriented gap closure.}
Round~9 reduced the unresolved part of the line
\[
|E(G)|=2|V(G)|-4
\]
for $43\le n\le 51$ to a single obstruction: any graph $G$ with
\[
|V(G)|=n,\qquad |E(G)|=2n-4,\qquad S(G)<\frac14
\]
must satisfy
\[
\mathcal C(G)=\{6,L\}
\qquad\text{for some }L\ge 13.
\]
In this round I eliminate that obstruction by proving a general density bound for graphs whose cycle-length set is $\{6,L\}$. This completes the whole previously unresolved interval $43\le n\le 51$.

\subsection*{2) Round-9 FOUNDATION USED}

I use the following vetted Round~9 results as black boxes.
\begin{itemize}
\item[(F1)] For $2\le a\le b$, the complete bipartite graph $K_{a,b}$ has
\[
\mathcal C(K_{a,b})=\{4,6,\dots,2a\},
\qquad
S(K_{a,b})=\sum_{j=2}^{a}\frac1{2j}=\frac12(H_a-1).
\]
In particular,
\[
S(K_{2,n-2})=\frac14.
\]

\item[(F2)] For every $7\le n\le 42$,
\[
\min\{S(G): |V(G)|=n,\ |E(G)|=2n-4\}=\frac14,
\]
attained by $K_{2,n-2}$.

\item[(F3)] If $43\le n\le 51$, $|V(G)|=n$, $|E(G)|=2n-4$, and $S(G)<1/4$, then
\[
\mathcal C(G)=\{6,L\}
\qquad\text{for some }L\ge 13.
\]

\item[(F4)] If $X$ is connected and every cycle of $X$ has length $6$ (or $X$ is acyclic), then
\[
|E(X)|\le \frac32\bigl(|V(X)|-1\bigr).
\]
\end{itemize}

\subsection*{3) NEW INSIGHT / TOOL (ROUND-10)}

The new ingredient is a cycle-started $H$-path growth bound for the two-length family $\{6,L\}$.
If a $2$-connected graph $B$ has no cycle lengths outside $\{6,L\}$ and contains an $L$-cycle, then every $H$-path used to enlarge an $L$-cycle inside $B$ must have length at least $3$.
Since an added path of length $s\ge 3$ contributes $s$ edges and only $s-1$ new vertices, each such step adds at most $3/2$ edges per new vertex. Starting from one $L$-cycle, this yields the block bound
\[
|E(B)|\le \frac32|V(B)|-\frac L2.
\]
Combined with the Round~9 estimate for the ``all cycles are $6$'' case, this rules out every graph with cycle-length set $\{6,L\}$ at density $2n-4$.

\subsection*{4) ATTACK PLAN (ROUND-10)}

The exact Round~9 gap was:
\[
43\le n\le 51,
\qquad
|V(G)|=n,
\qquad
|E(G)|=2n-4,
\qquad
\mathcal C(G)=\{6,L\}\text{ for some }L\ge 13.
\]
The plan is:
\begin{enumerate}
\item prove a self-contained density bound for every $2$-connected graph whose cycle lengths lie in $\{6,L\}$ and which contains an $L$-cycle;
\item extend it to connected graphs by induction on the block tree, using Foundation~(F4) on the blocks with no $L$-cycle;
\item extend it once more to disconnected graphs;
\item apply the result to the Round~9 obstruction.
\end{enumerate}
This turns the last obstruction into an outright contradiction, and therefore finishes the whole range $43\le n\le 51$.

\subsection*{5) WORK (ROUND-10)}

As before,
\[
\mathcal C(G):=\{\ell\ge 3:\text{$G$ contains a cycle of length $\ell$}\},
\qquad
S(G):=\sum_{\ell\in\mathcal C(G)}\frac1\ell.
\]

\subsubsection*{5.1. Basic lemmas on $H$-paths}

\textbf{Definition.}
If $H\subseteq G$ is a subgraph, an \emph{$H$-path} in $G$ is a path whose endpoints lie in $H$ and whose internal vertices lie outside $H$. We allow an $H$-path to have length $1$.

\medskip
\textbf{Lemma 5.1.}
Let $H$ be a proper $2$-connected subgraph of a $2$-connected graph $B$. Then $B$ contains an $H$-path.

\medskip
\textit{Proof.}
If $V(H)=V(B)$ but $E(H)\neq E(B)$, any edge of $E(B)\setminus E(H)$ is an $H$-path of length $1$.
So assume $V(H)\neq V(B)$ and let $K$ be a component of $B-V(H)$.
Because $B$ is connected, $K$ has a neighbor in $H$.
If it had only one such neighbor $x\in V(H)$, then deleting $x$ would separate $K$ from $H-x$, contradicting $2$-connectedness of $B$.
Hence $K$ has at least two neighbors in $H$.
A shortest path in $K$ between two vertices adjacent to $H$, together with its two incident edges into $H$, gives an $H$-path.
\hfill$\square$

\medskip
\textbf{Lemma 5.2.}
If $H$ is $2$-connected and $P$ is an $H$-path, then $H\cup P$ is $2$-connected.

\medskip
\textit{Proof.}
Let the endpoints of $P$ be $u,v\in V(H)$.
We show that deleting any vertex from $H\cup P$ leaves a connected graph.
If the deleted vertex is not on $P$, then $H$ minus that vertex is connected and still contains $u,v$, so it also contains the remaining path $P$.
If the deleted vertex is internal to $P$, then $P$ splits into two paths leading to $u$ and $v$, while $H$ minus that vertex is connected.
If the deleted vertex is $u$ (the case $v$ is symmetric), then $H-u$ is connected and contains $v$, so every remaining vertex of $P$ still reaches $H-u$ through the $v$-end.
Thus $(H\cup P)-z$ is connected for every vertex $z$, so $H\cup P$ is $2$-connected.
\hfill$\square$

\subsubsection*{5.2. One $2$-connected $\{6,L\}$-block}

\textbf{Theorem 5.3.}
Let $L\ge 13$, and let $B$ be a $2$-connected graph such that
\[
\mathcal C(B)\subseteq\{6,L\}
\qquad\text{and}\qquad
L\in\mathcal C(B).
\]
Then
\[
|E(B)|\le \frac32|V(B)|-\frac L2
= \frac32\bigl(|V(B)|-1\bigr)-\frac{L-3}{2}.
\]

\medskip
\textit{Proof.}
Choose an $L$-cycle $C\subseteq B$.
Let $H$ be a maximal $2$-connected subgraph of $B$ obtainable from $C$ by repeatedly adjoining $H$-paths.
We claim that $H=B$.

Assume for contradiction that $H\neq B$.
By Lemma~5.1, there is an $H$-path $P$ in $B$.
Let $s:=|E(P)|$ and let $u,v\in V(H)$ be its endpoints.
Since $H$ is $2$-connected, there exist two internally disjoint $u$--$v$ paths in $H$; their union contains a cycle $D\subseteq H$ through $u$ and $v$.
Let $m:=|D|\in\{6,L\}$, and let the two $u$--$v$ arcs of $D$ have lengths $d$ and $m-d$.
Then $P$ together with each arc forms a cycle of $B$, so
\[
s+d\in\{6,L\},
\qquad
s+m-d\in\{6,L\}.
\]
We show that $s\ge 3$.

If $m=6$ and $s\le 2$, then the two new cycle lengths have sum
\[
(s+d)+(s+6-d)=2s+6\in\{8,10\},
\]
whereas any sum of two numbers from $\{6,L\}$ is one of $12$, $L+6$, or $2L$, impossible for $L\ge 13$.

If $m=L$ and $s\le 2$, then the two new cycle lengths have sum
\[
(s+d)+(s+L-d)=2s+L\in\{L+2,L+4\},
\]
whereas again any sum of two numbers from $\{6,L\}$ is one of $12$, $L+6$, or $2L$.
So $s\ge 3$.

By Lemma~5.2, adjoining $P$ to $H$ gives a larger $2$-connected subgraph still obtainable from $C$ by adjoining $H$-paths, contradicting maximality of $H$.
Therefore $H=B$.

So $B$ can be built from the initial $L$-cycle by adjoining paths $P_1,\dots,P_t$ of lengths $s_1,\dots,s_t$, all at least $3$.
Each $P_i$ contributes $s_i$ new edges and $s_i-1$ new vertices, hence
\[
|E(B)|=L+\sum_{i=1}^t s_i,
\qquad
|V(B)|=L+\sum_{i=1}^t (s_i-1).
\]
Thus
\[
|E(B)|=|V(B)|+t.
\]
Also,
\[
|V(B)|-L=\sum_{i=1}^t (s_i-1)\ge 2t,
\qquad\text{so}\qquad
t\le \frac{|V(B)|-L}{2}.
\]
Substituting gives
\[
|E(B)|\le |V(B)|+\frac{|V(B)|-L}{2}=\frac32|V(B)|-\frac L2.
\]
\hfill$\square$

\subsubsection*{5.3. Connected and disconnected $\{6,L\}$-graphs}

\textbf{Proposition 5.4.}
Let $L\ge 13$, and let $G$ be a connected graph such that
\[
\mathcal C(G)\subseteq\{6,L\}
\qquad\text{and}\qquad
L\in\mathcal C(G).
\]
Then
\[
|E(G)|\le \frac32\bigl(|V(G)|-1\bigr)-\frac{L-3}{2}.
\]

\medskip
\textit{Proof.}
We induct on the number of blocks of $G$.
If $G$ has only one block, then because it contains an $L$-cycle it is not a bridge, hence it is $2$-connected, and the claim is exactly Theorem~5.3.

Assume now that $G$ has at least two blocks.
Choose a leaf block $B$ in the block tree of $G$, and let $H$ be the union of the other blocks.
Then $H$ is connected and
\[
|V(G)|=|V(B)|+|V(H)|-1,
\qquad
|E(G)|=|E(B)|+|E(H)|.
\]

If $B$ contains an $L$-cycle, then Theorem~5.3 gives
\[
|E(B)|\le \frac32\bigl(|V(B)|-1\bigr)-\frac{L-3}{2}.
\]
For $H$ we have the weaker bound
\[
|E(H)|\le \frac32\bigl(|V(H)|-1\bigr),
\]
either by the inductive hypothesis (if $H$ also contains an $L$-cycle) or by Foundation~(F4) (if it does not).
Summing gives the result.

If $B$ contains no $L$-cycle, then every cycle of $B$ has length $6$, or else $B\cong K_2$.
So Foundation~(F4) gives
\[
|E(B)|\le \frac32\bigl(|V(B)|-1\bigr).
\]
Since $G$ contains an $L$-cycle and $B$ does not, the remainder $H$ still contains an $L$-cycle, so the inductive hypothesis applies to $H$:
\[
|E(H)|\le \frac32\bigl(|V(H)|-1\bigr)-\frac{L-3}{2}.
\]
Again summing gives the result.
\hfill$\square$

\medskip
\textbf{Corollary 5.5.}
Let $L\ge 13$, and let $G$ be an arbitrary graph with
\[
\mathcal C(G)=\{6,L\}.
\]
If $c(G)$ denotes the number of connected components, then
\[
|E(G)|\le \frac32\bigl(|V(G)|-c(G)\bigr)-\frac{L-3}{2}.
\]
In particular,
\[
|E(G)|<2|V(G)|-4.
\]

\medskip
\textit{Proof.}
Choose one component $G_0$ containing an $L$-cycle.
By Proposition~5.4,
\[
|E(G_0)|\le \frac32\bigl(|V(G_0)|-1\bigr)-\frac{L-3}{2}.
\]
Every other component $G_i$ satisfies the weaker estimate
\[
|E(G_i)|\le \frac32\bigl(|V(G_i)|-1\bigr),
\]
by Proposition~5.4 if it also contains an $L$-cycle, and by Foundation~(F4) otherwise.
Summing over all components gives
\[
|E(G)|\le \frac32\sum_i\bigl(|V(G_i)|-1\bigr)-\frac{L-3}{2}
= \frac32\bigl(|V(G)|-c(G)\bigr)-\frac{L-3}{2}.
\]
Since $c(G)\ge 1$,
\[
|E(G)|\le \frac32\bigl(|V(G)|-1\bigr)-\frac{L-3}{2}.
\]
Finally,
\[
2|V(G)|-4-\left(\frac32\bigl(|V(G)|-1\bigr)-\frac{L-3}{2}\right)
=\frac{|V(G)|+L-8}{2}>0,
\]
so $|E(G)|<2|V(G)|-4$.
\hfill$\square$

\subsubsection*{5.4. Completion of the previously unresolved interval}

\textbf{Theorem 5.6.}
For every integer $43\le n\le 51$,
\[
\min\{S(G): |V(G)|=n,\ |E(G)|=2n-4\}=\frac14.
\]
The minimum is attained by the complete bipartite graph $K_{2,n-2}$.

\medskip
\textit{Proof.}
By Foundation~(F1), $S(K_{2,n-2})=1/4$.
So it remains to prove that every graph $G$ on $n$ vertices with $2n-4$ edges satisfies $S(G)\ge 1/4$.
Assume for contradiction that $S(G)<1/4$.
Then Foundation~(F3) gives
\[
\mathcal C(G)=\{6,L\}
\qquad\text{for some }L\ge 13.
\]
But Corollary~5.5 then implies
\[
|E(G)|<2|V(G)|-4,
\]
contradicting $|E(G)|=2n-4$.
Therefore $S(G)\ge 1/4$, and the example $K_{2,n-2}$ shows that the minimum is exactly $1/4$.
\hfill$\square$

\medskip
\textbf{Corollary 5.7.}
For every integer $7\le n\le 51$,
\[
\min\{S(G): |V(G)|=n,\ |E(G)|=2n-4\}=\frac14,
\]
attained by $K_{2,n-2}$.

\medskip
\textit{Proof.}
For $7\le n\le 42$ this is Foundation~(F2). For $43\le n\le 51$ it is Theorem~5.6.
\hfill$\square$

\subsection*{6) ADVERSARIAL VERIFICATION}

\begin{itemize}
\item \textbf{$H$-path existence.} Lemma~5.1 explicitly handles both possibilities: either $H$ is proper because of missing vertices or because of missing edges.

\item \textbf{Preservation of $2$-connectedness.} Lemma~5.2 is proved directly by checking connectivity after deleting an arbitrary vertex; no external ear-decomposition theorem is used.

\item \textbf{The key inequality $|P|\ge 3$.} In Theorem~5.3 the contradiction for $|P|=1$ or $2$ depends only on the two arc lengths of one cycle through the endpoints and the fact that both resulting cycle lengths must lie in $\{6,L\}$. No stronger classification claim is required.

\item \textbf{Block induction.} In Proposition~5.4 the identities
\[
|V(G)|=|V(B)|+|V(H)|-1,
\qquad
|E(G)|=|E(B)|+|E(H)|
\]
are valid because a leaf block meets the rest of the graph in exactly one cutvertex. The case split according to whether $B$ contains an $L$-cycle is exhaustive.

\item \textbf{Interaction with Round~9.} Round~9 reduced any putative counterexample in $43\le n\le 51$ to the form $\mathcal C(G)=\{6,L\}$. Corollary~5.5 shows that no graph of that form can reach $2n-4$ edges, so the new argument genuinely closes the last Round~9 obstruction.
\end{itemize}

\subsection*{7) FINAL (EXACTLY ONE)}

\textbf{UNRESOLVED (BUT STRICTLY ADVANCED).}

Round~10 completes the entire unresolved interval left after Round~9. It proves that for every $43\le n\le 51$,
\[
\min\{S(G): |V(G)|=n,\ |E(G)|=2n-4\}=\frac14,
\]
with extremal graph $K_{2,n-2}$.
Hence the complete-bipartite minimiser statement on the line $m=2n-4$ is now fully proved for
\[
7\le n\le 51.
\]
The global Erd\H{o}s--Hajnal minimiser question for arbitrary $(n,m)$ remains open.

\subsection*{8) COMPLETION ESTIMATE (MANDATORY)}

\textbf{COMPLETION: 99\%}

\subsection*{9) REFERENCES}

No new external references were used in the new proofs of this round. The only imported inputs are the vetted Round~9 results listed in \S2.
